\SourceLeaf{065}%
tuendos inter se comparare licebit: optimaque combinatio ea erit, in qua hic error medius quam minimus evadit. Qui quum fiat
\[
=\BareRoot{(LLmm+L'L'm'm'+L''L''m''m''+\text{etc.})}
\]
in id erit incumbendum, ut aggregatum \(LLmm+L'L'm'm'+L''L''m''m''+\text{etc.}\) nanciscatur valorem minimum.

\Article{4}
Quum varietas infinita functionum \(U\), quae secundum conditionem in art. praec. enunciatam ipsius \(u\) vice fungi possunt, eatenus tantum hic consideranda veniat, quatenus diversa systemata valorum coëfficientium \(L,L',L''\) etc. inde sequuntur, indagare oportebit ante omnia nexum, qui inter cuncta systemata admissibilia locum habere debet. Designemus valores determinatos quotientium differentialium partialium
\[
\begin{array}{ccc l}
\dfrac{\mathrm dX}{\mathrm dv},&\dfrac{\mathrm dX}{\mathrm dv'},&\dfrac{\mathrm dX}{\mathrm dv''}&\text{etc.}\\[5pt]
\dfrac{\mathrm dY}{\mathrm dv},&\dfrac{\mathrm dY}{\mathrm dv'},&\dfrac{\mathrm dY}{\mathrm dv''}&\text{etc.}\\[5pt]
\dfrac{\mathrm dZ}{\mathrm dv},&\dfrac{\mathrm dZ}{\mathrm dv'},&\dfrac{\mathrm dZ}{\mathrm dv''}&\text{etc. etc.}
\end{array}
\]
quos obtinent, si ipsis \(v,v',v''\) etc. valores veri tribuuntur, resp. per
\[
\begin{array}{ccc l}
a,&a',&a''&\text{etc.}\\
b,&b',&b''&\text{etc.}\\
c,&c',&c''&\text{etc. etc.}
\end{array}
\]
patetque, si ipsis \(v,v',v''\) etc. accedere concipiantur talia incrementa \(\mathrm dv,\mathrm dv',\mathrm dv''\) etc., per quae \(X,Y,Z\) etc. non mutentur, adeoque singulae maneant \(=0\), i. e. satisfacientia aequationibus
\[
\begin{aligned}
0&=a\mathrm dv+a'\mathrm dv'+a''\mathrm dv''+\text{etc.}\\
0&=b\mathrm dv+b'\mathrm dv'+b''\mathrm dv''+\text{etc.}\\
0&=c\mathrm dv+c'\mathrm dv'+c''\mathrm dv''+\text{etc. etc.}
\end{aligned}
\]
etiam \(u-U\) non mutari debere, adeoque fieri
\[
0=(l-L)\mathrm dv+(l'-L')\mathrm dv'+(l''-L'')\mathrm dv''+\text{etc.}
\]
\SourceLeaf{066}%
Hinc facile concluditur, coëfficientes \(L,L',L''\) etc. contentos esse debere sub formulis talibus
\[
\begin{aligned}
L&=l+ax+by+cz+\text{etc.}\\
L'&=l'+a'x+b'y+c'z+\text{etc.}\\
L''&=l''+a''x+b''y+c''z+\text{etc.}
\end{aligned}
\]
etc., denotantibus \(x,y,z\) etc. multiplicatores determinatos. Vice versa patet, si systema multiplicatorum determinatorum \(x,y,z\) etc. ad lubitum assumatur, semper assignari posse functionem \(U\) talem, cui valores ipsorum \(L,L',L''\) etc. his aequationibus conformes respondeant, et quae pro conditione in art. praec. enunciata ipsius \(u\) vice fungi possit: quin adeo hoc infinitis modis diversis effici posse. Modus simplicissimus erit statuere \(U=u+xX+yY+zZ+\text{etc.}\); generalius statuere licet \(U=u+xX+yY+zZ+\text{etc.}+u'\), denotante \(u'\) talem functionem indeterminatarum \(v,v',v''\) etc., quae semper evanescit pro \(X=0,\ Y=0,\ Z=0\) etc., et cuius valor in casu determinato de quo agitur fit maximus vel minimus. Sed ad institutum nostrum nulla hinc oritur differentia.

\Article{5}
Facile iam erit, multiplicatoribus \(x,y,z\) etc. valores tales tribuere, ut aggregatum
\[
LLmm+L'L'm'm'+L''L''m''m''+\text{etc.}
\]
assequatur valorem minimum. Manifesto ad hunc finem haud opus est cognitione errorum mediorum \(m,m',m''\) etc. absoluta, sed sufficit ratio, quam inter se tenent. Introducemus itaque ipsorum loco pondera observationum \(p,p',p''\) etc., i. e. numeros quadratis \(mm,m'm',m''m''\) etc. reciproce proportionales, pondere alicuius observationis ad lubitum pro unitate accepto. Quantitates \(x,y,z\) etc. itaque sic determinari debebunt, ut polynomium indefinitum
\[
\begin{gathered}
\frac{(ax+by+cz+\text{etc.}+l)^2}{p}+\frac{(a'x+b'y+c'z+\text{etc.}+l')^2}{p'}\\[4pt]
{}+\frac{(a''x+b''y+c''z+\text{etc.}+l'')^2}{p''}+\text{etc.}
\end{gathered}
\]
nanciscatur valorem minimum, quod fieri supponemus per valores \textit{determinatos} \(x^0,y^0,z^0\) etc.

Introducendo denotationes sequentes
\SourceLeaf{067}
\[
\begin{aligned}
\frac{aa}{p}+\frac{a'a'}{p'}+\frac{a''a''}{p''}+\text{etc.}&=[aa]\\[4pt]
\frac{ab}{p}+\frac{a'b'}{p'}+\frac{a''b''}{p''}+\text{etc.}&=[ab]\\[4pt]
\frac{ac}{p}+\frac{a'c'}{p'}+\frac{a''c''}{p''}+\text{etc.}&=[ac]\\[4pt]
\frac{bb}{p}+\frac{b'b'}{p'}+\frac{b''b''}{p''}+\text{etc.}&=[bb]\\[4pt]
\frac{bc}{p}+\frac{b'c'}{p'}+\frac{b''c''}{p''}+\text{etc.}&=[bc]\\[4pt]
\frac{cc}{p}+\frac{c'c'}{p'}+\frac{c''c''}{p''}+\text{etc.}&=[cc]
\end{aligned}
\]
\begin{center}etc. nec non\end{center}
\[
\begin{aligned}
\frac{al}{p}+\frac{a'l'}{p'}+\frac{a''l''}{p''}+\text{etc.}&=[al]\\[4pt]
\frac{bl}{p}+\frac{b'l'}{p'}+\frac{b''l''}{p''}+\text{etc.}&=[bl]\\[4pt]
\frac{cl}{p}+\frac{c'l'}{p'}+\frac{c''l''}{p''}+\text{etc.}&=[cl]
\end{aligned}
\]
\begin{center}etc.\end{center}
manifesto conditio minimi requirit, ut fiat
\[
\left.\begin{aligned}
0&=[aa]x^0+[ab]y^0+[ac]z^0+\text{etc.}+[al]\\
0&=[ab]x^0+[bb]y^0+[bc]z^0+\text{etc.}+[bl]\\
0&=[ac]x^0+[bc]y^0+[cc]z^0+\text{etc.}+[cl]\\
&\text{etc.}
\end{aligned}\right\}\ (1)
\]
Postquam quantitates \(x^0,y^0,z^0\) etc. per eliminationem hinc derivatae sunt, statuetur
\[
\left.\begin{aligned}
ax^0+by^0+cz^0+\text{etc.}+l&=L\\
a'x^0+b'y^0+c'z^0+\text{etc.}+l'&=L'\\
a''x^0+b''y^0+c''z^0+\text{etc.}+l''&=L''\\
\text{etc.}\quad&
\end{aligned}\right\}\ (2)
\]
His ita factis, functio quantitatum \(v,v',v''\) etc. ea ad determinationem incognitae nostrae maxime idonea minimaeque incertitudini obnoxia erit, cuius quotientes differentiales partiales in casu determinato de quo agitur habent valores \(L,L',L''\) etc. resp., pondusque huius determinationis, quod per \(P\) denotabimus, erit
